\begin{center}
    \begin{asy}
        import olympiad;
        import geometry;
        // Set the size of the diagram
        size(250);

        // Define the vertices of triangle ABC
        pair A, B, C;
        B = (0,0);
        A = (1,6);
        C = (8,1);

        // Define the direction of Line L (perpendicular to AB)
        pair L_dir = rotate(90)*(A-B);
        // Draw the line L extending in both directions from B
        path L = (B - 7*unit(L_dir)) -- (B + 10*unit(L_dir));

        // The line from A perpendicular to BC
        pair A_perp_foot = foot(A, B, C);
        path line_AD = A--A_perp_foot;

        // Point D is the intersection of L and the perpendicular from A to BC
        pair D = extension(B, B + L_dir, A, A_perp_foot);

        // The perpendicular bisector of BC
        pair M = (B+C)/2; // Midpoint of BC
        pair perp_bisector_dir = rotate(90)*(C-B);
        path perp_bisector = M -- M + 5*unit(perp_bisector_dir);

        // Point P is the intersection of L and the perpendicular bisector of BC
        pair P = extension(B, D, M, M + perp_bisector_dir);

        // Point E is the foot of the perpendicular from D to AC
        pair E = foot(D, A, C);

        // Point H is the intersection of BE and the perpendicular bisector of BC
        pair N = extension(B, E, M, M + perp_bisector_dir);

        // T intersection of AD and BC
        pair T = intersectionpoint(line(B,C),line(A,D));

        // --- Drawing Commands ---

        // Draw the main triangle ABC
        draw(A--B--C--A);

        // Draw the construction lines
        draw(A--D, red);
        draw(M--P--N, darkgreen);

        // Draw the triangle BPE, which is to be proved isosceles
        draw(B--P--E--P--C, magenta);
        draw(B--E);
        draw(N--C);

        // Add dots and labels for all key points
        dot("$A$", A, NW);
        dot("$B$", B, SW);
        dot("$C$", C, SE);
        dot("$D$", D, SW);
        dot("$P$", P, SW);
        dot("$E$", E, NE);
        dot("$M$", M, SW);
        dot("$N$", N, NW);
        dot("$T$", T, NE);

        // Add right angle markers to show perpendicular lines
        rightanglemark(A, B, D, 5);
        rightanglemark(A, A_perp_foot, C, 5);
        rightanglemark(D, E, A, 5);
        rightanglemark(P, M, B, 5);
    \end{asy}
\end{center}